\section{Methods}
\label{sec:methods}

\subsection{Overview}
\label{sec:overview}

Our pipeline consumes a human motion-capture clip and GMR's own per-frame
retarget of it, and produces a contact-corrected, smoothed, floor-safe
trajectory for the Unitree G1 humanoid. Every stage after GMR's own solve is
ours; GMR itself (its tracking cost, its joint-limit handling, its solver)
is used completely unmodified. We present two instantiations of the
correction layer that sits on top of GMR's output: a \textbf{multi-stage}
pipeline built from a sequence of per-frame corrections
(Section~\ref{sec:multistage}), and a \textbf{global} pipeline that replaces
that sequence with a single whole-trajectory optimization
(Section~\ref{sec:global}). Both share the same canonical-human contact
labeling (Section~\ref{sec:canonical}), the same evaluation protocol
(Section~\ref{sec:grounding}), and the same final grounding step.
Section~\ref{sec:comparing} compares them.

Throughout this paper, ``contact-aware'' and ``contact-first'' refer to
\textit{geometric} contact handling: per-frame detection of which effector
is in contact (Section~\ref{sec:canonical}), and enforcement of
floor-clearance, held-position, and self-collision constraints derived from
that detection. Neither pipeline models contact forces, friction, or joint
torque --- the dynamics of contact are outside this paper's scope by design
(Sections~\ref{sec:introduction} and~\ref{sec:related-work}), the same
scope every kinematic retargeter we compare against shares.
\cite{ayusawa2017morphing}'s simplified zero-moment-point check is the
closest any of them come to a dynamics proxy, and even it assumes
quasi-static balance rather than modeling contact forces directly. The
planned policy-training evaluation (Section~\ref{sec:status-next-steps}) is
the step that tests whether a geometrically contact-correct reference is
also dynamically trackable; this paper's contribution stops at the
kinematic guarantee.

\subsection{Canonical Human and Contact Detection}
\label{sec:canonical}

Raw BVH motion capture (LAFAN1, 30~fps) is mapped to a robot-agnostic
canonical skeleton: per-frame landmark positions and orientations derived
from bone geometry, with hand orientation taken directly from GMR's own
raw-bone-rotation signal to preserve wrist twist. Contact is a per-frame
detection gate, not a correction: an effector is flagged ``in contact'' only
when it is simultaneously below a height threshold and moving slower than
0.4~m/s (7~cm for feet, 8~cm for hands) --- height alone is insufficient,
since a fast-swinging limb can pass close to the floor without being
planted. Crossing the threshold marks a frame as a candidate for anchoring
by a later stage; it does not move anything by itself. This detector covers
feet and hands and is shared verbatim by every downstream stage that needs
to know which frames are in contact.

For corpus-level evaluation, clips are additionally classified into a
floor-contact class and a locomotion class using a broader six-region human
contact labeler (feet, hands, knees, elbows, pelvis, torso; 5/8/15~cm
respectively), since a feet-only or hip-height signal undercounts floor
contact on clips where support is briefly taken through a hand or knee.

\subsection{Multi-Stage Pipeline}
\label{sec:multistage}

\noindent\textbf{Floor and self-collision clamp.} For each frame, a
two-phase sequential damped-least-squares correction runs per limb chain
(hip$\to$knee$\to$ankle, shoulder$\to$elbow$\to$wrist), proximal-to-distal.
Phase~1 resolves floor violations and pins held effectors to their contact
target; Phase~2, run strictly afterward rather than jointly with Phase~1,
resolves self-collision using the same contact-normal and
relative-Jacobian formulation as the collision term in the global pipeline
(Section~\ref{sec:global}).

\noindent\textbf{Held-aware smoothing.} A closed-form per-joint tridiagonal
smoother removes joint-velocity spikes. Held-effector degrees of freedom are
locked to their input value at high weight during contact, ramped in and
out over five frames, so smoothing does not erode a contact correction the
clamp stage already made. The floor/self-collision clamp is re-applied once
more after smoothing, since geometry-blind smoothing can reintroduce a
violation at a non-held frame.

\noindent\textbf{Rate-limited correction.} A temporal trust region on the
applied correction bounds how fast it can change from one frame to the
next, independent of the pose itself. This mitigates branch-flip artifacts
(discrete jumps between distinct, equally-valid per-frame solutions on a
redundant kinematic chain --- Section~\ref{sec:disc-tradeoff}) without
addressing their cause.

\noindent\textbf{Grounding.} See Section~\ref{sec:grounding} below (shared
with the global pipeline).

\subsection{Global Pipeline}
\label{sec:global}

The global pipeline replaces the clamp / smoothing / rate-limit sequence in
Section~\ref{sec:multistage} with one whole-trajectory optimization. GMR's
raw per-frame output is consumed directly as input.

\noindent\textbf{Floor and self-collision correction.} A sparse QP over the
actuated joint trajectory of the entire clip, solved by sequential convex
approximation: at each outer iteration, floor and self-collision
constraints are linearized into soft, slack-penalized rows (never
primal-infeasible), a trust region bounds the step, and a smoothness term
penalizes the change in the SCA correction between adjacent frames, in the
same objective as the floor, collision, and contact-tracking terms. Held
effectors are anchored to a per-region target: contact timing comes from
the canonical human's own labels (Section~\ref{sec:canonical}), a stillness
sub-check on the robot's own trajectory splits each labeled contact
interval into stationary sub-runs (anchored to a single position, ramped
in/out over five frames) and repositioning sub-frames (lightly regularized
toward their own per-frame position), and the anchor's target height is a
fixed per-effector geometric constant (the robot's own origin-to-support-point
offset at a neutral pose), not a value computed from the possibly-unreliable
input trajectory. Floor penetration is intentionally excluded from this
stage's own convergence scoring, since its solve variable is the robot's
joints only --- a pelvis-level violation on a floor-contact pose is outside
what this stage can correct, and grounding (Section~\ref{sec:grounding})
resolves it regardless of this stage's outcome; scoring against it would
only reject otherwise-improving self-collision corrections. A
cross-iteration keep-best selection guarantees the result is never worse
than GMR's own raw output on this stage's own objective.

\noindent\textbf{Smoothing and grounding.} The corrected trajectory passes
through the same held-aware smoothing stage used in
Section~\ref{sec:multistage}, then grounding (Section~\ref{sec:grounding}).
No per-frame re-clamp or rate limiter is applied --- the whole-trajectory
formulation's own smoothness term is the temporal-coupling mechanism, in
place of the multi-stage pipeline's rate limiter.

Because its smoothness term shares an objective with its floor, collision,
and tracking terms, a solution that alternates between distinct null-space
branches on adjacent frames costs the objective directly; the global
pipeline cannot produce a branch-flip by construction, independent of any
per-frame tuning.

\subsection{Grounding and Evaluation}
\label{sec:grounding}

\noindent\textbf{Grounding.} A final, deterministic pass computes the
per-frame vertical shift required to clear the floor, applies a causal
max-filter and Gaussian smoothing to that requirement, and takes the
pointwise maximum against the raw requirement. This guarantees the shifted
trajectory has zero floor penetration by construction. Unlike GMR's own
height-fix --- one rigid per-clip shift --- this is computed per frame, so
it does not force a single Z-offset to serve both a clip's standing phase
and its lying phase.

\noindent\textbf{Metrics.} No ground-truth robot motion exists, so every
metric is computed from the output trajectory and the robot's own vetted
collision geometry: floor penetration and self-collision incidence
(mesh-exact lowest point vs.\ $z=0$, k-hop-filtered for anatomically
adjacent pairs); worst float (height of a held effector's support point
above the floor); tracking fidelity (FK'd body position/orientation vs.\
GMR's own scaled human target); jerk, peak joint velocity, and
velocity-spike count; and skate (horizontal drift of a held contact point
from its own contact-onset position).

\noindent\textbf{Joint contact-correctness.} A held frame passes only if
its support point is simultaneously within a tight band of the floor and
the whole-body mesh clears the floor everywhere at that same frame. Both
conditions are necessary: a metric that checks either alone is satisfiable
by a per-clip constant Z-shift with no notion of per-frame contact at all,
since a rigid shift can be tuned to land in the right place at a clip's
labeled held frames while leaving the rest of the clip penetrating or
floating arbitrarily. Requiring both simultaneously is not satisfiable by
any single rigid shift on a multi-phase clip
(Section~\ref{sec:results-oracle}), which is why we report this joint
metric as the primary measure of contact quality rather than penetration or
float in isolation.

\subsection{Comparing the Two Pipelines}
\label{sec:comparing}

The multi-stage pipeline (Section~\ref{sec:multistage}) benefits from a
large total per-frame solve budget --- a full iterative local correction at
every frame --- at the cost of no explicit temporal coupling between
frames, which manifests as branch-flip artifacts and their associated
joint-velocity cost. The global pipeline (Section~\ref{sec:global}) removes
that artifact by construction but currently allocates a smaller total solve
budget (a bounded number of whole-clip outer iterations) to contact
precision. Section~\ref{sec:results} reports both.
